\section*{NeurIPS Paper Checklist}

\begin{enumerate}

\item {\bf Claims}
    \item[] Question: Do the main claims made in the abstract and introduction accurately reflect the paper's contributions and scope?
    \item[] Answer: \answerYes{}
    \item[] Justification: The abstract distinguishes exact invariance of fixed independently encoded elements from empirical behavior of contextual representations, identifies pair-classification results as diagnostics, and avoids causal attribution. Results invalidated by implementation or protocol changes are marked pending rerun.

\item {\bf Limitations}
    \item[] Question: Does the paper discuss the limitations of the work performed by the authors?
    \item[] Answer: \answerYes{}
    \item[] Justification: Section~5 discusses DTW scalability, probe multiplicity, frame-count confounds, dependent-pair evaluation, the small selected SSv2 split, and limits on cross-dataset generalization.

\item {\bf Theory assumptions and proofs}
    \item[] Question: For each theoretical result, does the paper provide the full set of assumptions and a complete (and correct) proof?
    \item[] Answer: \answerNA{}
    \item[] Justification: The paper makes no new formal theorem claim. The stated permutation invariance of a symmetric function over fixed inputs follows from its definition, and its limited assumptions are made explicit.

    \item {\bf Experimental result reproducibility}
    \item[] Question: Does the paper fully disclose all the information needed to reproduce the main experimental results of the paper to the extent that it affects the main claims and/or conclusions of the paper (regardless of whether the code and data are provided or not)?
    \item[] Answer: \answerNo{}
    \item[] Justification: The code and guide specify the current protocols, but corrected scramble, EPIC residual, paired-bootstrap, and directed-reranking results are pending. Exact feature caches and a locked environment are not included. This answer should be revisited after the documented cluster jobs complete.

\item {\bf Open access to data and code}
    \item[] Question: Does the paper provide open access to the data and code, with sufficient instructions to faithfully reproduce the main experimental results, as described in supplemental material?
    \item[] Answer: \answerNo{}
    \item[] Justification: Code and instructions are provided, but several datasets require separate licenses or access agreements (notably Honda HDD), model weights have their own terms, and feature caches are not redistributed. The guide documents these requirements.

\item {\bf Experimental setting/details}
    \item[] Question: Does the paper specify all the training and test details (e.g., data splits, hyperparameters, how they were chosen, type of optimizer) necessary to understand the results?
    \item[] Answer: \answerYes{}
    \item[] Justification: No training is performed; all experiments use frozen pretrained models. Evaluation details (pair construction, DBSCAN $\epsilon$, negative-sampling ratios, bootstrap parameters, DTW $\alpha$ values, frame counts) are specified in Sections~2--3 and the appendix.

\item {\bf Experiment statistical significance}
    \item[] Question: Does the paper report error bars suitably and correctly defined or other appropriate information about the statistical significance of the experiments?
    \item[] Answer: \answerNo{}
    \item[] Justification: Existing pair-resampling intervals are now labeled descriptive because pairs share clips. HDD and nuScenes require paired cluster-bootstrap contrasts, and analogous independent-unit analyses are not yet available for every benchmark. Prompt standard deviations measure prompt sensitivity, not sampling uncertainty.

\item {\bf Experiments compute resources}
    \item[] Question: For each experiment, does the paper provide sufficient information on the computer resources (type of compute workers, memory, time of execution) needed to reproduce the experiments?
    \item[] Answer: \answerNo{}
    \item[] Justification: The paper reports the GPU class and approximate microbenchmarks, but it does not yet provide measured wall time and peak memory for every experiment. The supplied SLURM jobs provide requested resources for the corrected runs.

\item {\bf Code of ethics}
    \item[] Question: Does the research conducted in the paper conform, in every respect, with the NeurIPS Code of Ethics \url{https://neurips.cc/public/EthicsGuidelines}?
    \item[] Answer: \answerYes{}
    \item[] Justification: The research uses only publicly available datasets and pretrained models. The reversal attack analysis in Section~3.1 is diagnostic (identifying a vulnerability in existing systems) rather than offensive.

\item {\bf Broader impacts}
    \item[] Question: Does the paper discuss both potential positive societal impacts and negative societal impacts of the work performed?
    \item[] Answer: \answerYes{}
    \item[] Justification: The paper identifies a concrete Trust \& Safety vulnerability (reversal attack bypassing VCD filters) and diagnoses temporal blindness in video retrieval systems deployed at scale. The diagnostic framework itself poses minimal misuse risk as it characterizes existing model behavior rather than introducing new attack vectors.

\item {\bf Safeguards}
    \item[] Question: Does the paper describe safeguards that have been put in place for responsible release of data or models that have a high risk for misuse (e.g., pre-trained language models, image generators, or scraped datasets)?
    \item[] Answer: \answerNA{}
    \item[] Justification: The paper does not release new models or datasets. It evaluates existing publicly available models on existing public benchmarks.

\item {\bf Licenses for existing assets}
    \item[] Question: Are the creators or original owners of assets (e.g., code, data, models), used in the paper, properly credited and are the license and terms of use explicitly mentioned and properly respected?
    \item[] Answer: \answerYes{}
    \item[] Justification: All datasets and models are cited with their original publications. Appendix~\ref{sec:licenses} lists the governing license for each asset (datasets: VCDB research-only, Honda HDD HRI-USA Data Use Agreement, nuScenes CC BY-NC-SA 4.0, EPIC-Kitchens-100 CC BY-NC 4.0, Nymeria CC BY-NC 4.0, MUVR research-only; model weights listed alongside). All use is non-commercial research, consistent with the most restrictive license.

\item {\bf New assets}
    \item[] Question: Are new assets introduced in the paper well documented and is the documentation provided alongside the assets?
    \item[] Answer: \answerNA{}
    \item[] Justification: The paper does not release new datasets or models.

\item {\bf Crowdsourcing and research with human subjects}
    \item[] Question: For crowdsourcing experiments and research with human subjects, does the paper include the full text of instructions given to participants and screenshots, if applicable, as well as details about compensation (if any)?
    \item[] Answer: \answerNA{}
    \item[] Justification: No crowdsourcing or human subjects research was conducted.

\item {\bf Institutional review board (IRB) approvals or equivalent for research with human subjects}
    \item[] Question: Does the paper describe potential risks incurred by study participants, whether such risks were disclosed to the subjects, and whether Institutional Review Board (IRB) approvals (or an equivalent approval/review based on the requirements of your country or institution) were obtained?
    \item[] Answer: \answerNA{}
    \item[] Justification: No human subjects research was conducted.

\item {\bf Declaration of LLM usage}
    \item[] Question: Does the paper describe the usage of LLMs if it is an important, original, or non-standard component of the core methods in this research? Note that if the LLM is used only for writing, editing, or formatting purposes and does \emph{not} impact the core methodology, scientific rigor, or originality of the research, declaration is not required.
    \item[] Answer: \answerYes{}
    \item[] Justification: Three Video-LLMs (Qwen3-VL, Gemma~4, LLaVA-Video) are evaluated as subjects of the diagnostic study. Their usage is fully described in Sections~2--3. LLMs were also used to assist with writing and code development.

\end{enumerate}
